% ========================================
% CHƯƠNG 6: KẾT LUẬN VÀ HƯỚNG PHÁT TRIỂN
% ========================================
\chapter{KẾT LUẬN VÀ HƯỚNG PHÁT TRIỂN}
\label{ch:ketluan}

% ----- TÓM TẮT CHƯƠNG -----
\vspace{0.5cm}
\noindent\textit{Chương này tổng kết lại toàn bộ các kết quả nghiên cứu,
thiết kế và triển khai thực tế đã đạt được trong phạm vi của đồ án tốt
nghiệp. Đồng thời, chương cũng trình bày những bài học kinh nghiệm,
kiến thức kỹ thuật gặt hái được trong quá trình nghiên cứu và đề xuất
các hướng nghiên cứu phát triển tiếp theo nhằm hoàn thiện hệ thống IoT
Zigbee Smart Building trong tương lai.}
\vspace{0.5cm}

% ========================================
% 6.1 KẾ LƯẬN
% ========================================
\section{Kết luận chung}
\label{sec:ch6-conclusion}

Sau thời gian nghiên cứu và thực hiện, đồ án tốt nghiệp đã hoàn thành
đầy đủ các mục tiêu đề ra ban đầu, xây dựng thành công một hệ thống IoT
Smart Building đồng bộ từ tầng vật lý đến tầng ứng dụng của người dùng.
Các kết quả cụ thể đạt được bao gồm:

\textbf{Thiết kế và triển khai mạng Zigbee 3.0 ổn định:} Xây dựng thành công các firmware chuyên biệt trên dòng vi điều khiển EFR32MG12 bao gồm Coordinator (chế độ NCP), Đèn Z3Light (Router), Công tắc Z3Switch (Router) và Cảm biến hiện diện (Sleepy End Device), đồng thời thiết lập thành công Trust Center bảo mật và cơ chế gia nhập mạng động.

\textbf{Xây dựng Gateway nhúng hiệu năng cao:} Triển khai ứng dụng Gateway viết bằng ngôn ngữ C kết nối trực tiếp với thư viện Paho MQTT C để loại bỏ hoàn toàn các lớp trung gian Unix socket/IPC cũ. Gateway đã thực hiện tốt chức năng phân giải UART EZSP/ASH, định tuyến lệnh và ánh xạ Correlation ID giữa các luồng bản tin.

\textbf{Phát triển Cloud Backend và CSDL đám mây:} Xây dựng hệ thống API RESTful và MQTT Client bất đồng bộ sử dụng framework FastAPI kết nối PostgreSQL, triển khai Worker quét dọn Timeout lệnh ở nền và đóng gói toàn bộ hệ thống bằng Docker Compose trên AWS EC2.

\textbf{Phát triển ứng dụng di động Flutter:} Thiết kế giao diện ứng dụng di động trực quan cho phép giám sát trạng thái thiết bị gần thời gian thực, gửi lệnh bật/tắt Đèn và theo dõi lịch sử sự kiện hệ thống.

Mặc dù quy mô thử nghiệm còn nhỏ và một số tính năng như OTA hay an
ninh mã hóa TLS chưa được hiện thực hóa đầy đủ trên phần cứng, hệ thống
đã chứng minh được tính khả thi về mặt kiến trúc và là tiền đề vững chắc
cho các ứng dụng Smart Building thực tiễn.

% ========================================
% 6.2 KIẾN THỨC VÀ KỸ NĂNG ĐẠT ĐƯỢC
% ========================================
\section{Kiến thức và kỹ năng đạt được}
\label{sec:ch6-skills-learned}

Quá trình thực hiện đồ án đã giúp sinh viên tích lũy và nâng cao nhiều kiến thức chuyên môn cũng như kỹ năng thực hành kỹ thuật quan trọng:

\textbf{Kiến thức lý thuyết:} Hiểu rõ cấu trúc phân lớp của giao thức vô tuyến Zigbee 3.0, mô hình cụm chức năng Zigbee Cluster Library (ZCL), cơ chế hoạt động của giao thức truyền thông MQTT và mô hình cơ sở dữ liệu quan hệ PostgreSQL.

\textbf{Kỹ năng lập trình nhúng:} Thành thạo việc cấu hình, lập trình firmware trên IDE Simplicity Studio v5 của Silicon Labs, làm việc với ngăn xếp EmberZNet stack và lập trình hệ thống bằng C trên Linux.

\textbf{Kỹ năng Cloud và DevOps:} Tiếp cận với các công nghệ xây dựng API bất đồng bộ với FastAPI, quản lý và vận hành CSDL qua SQLAlchemy, đóng gói container Docker và quy trình triển khai máy chủ thực tế trên nền tảng điện toán đám mây AWS EC2.

\textbf{Kỹ năng nghiên cứu và giải quyết vấn đề:} Rèn luyện tư duy phân tích thiết kế hệ thống có tính hệ thống, phương pháp kiểm thử thực tế dựa trên logs, khắc phục các lỗi phát sinh thực tế như CRLF dòng lệnh, xung đột cổng dịch vụ hay phân quyền ACL.

% ========================================
% 6.3 HƯỚNG PHÁT TRIỂN TIẾP THEO
% ========================================
\section{Hướng phát triển tiếp theo}
\label{sec:ch6-future-work}

Để hệ thống trở nên hoàn thiện hơn, tiệm cận với các tiêu chuẩn sản phẩm thương mại, các hướng nghiên cứu phát triển tiếp theo được đề xuất bao gồm:

\subsection{Hoàn thiện phần cứng và mạng thiết bị}
\label{sec:ch6-fw-future}

\textbf{Mở rộng loại thiết bị:} Nghiên cứu phát triển thêm các loại node thiết bị mới như thiết bị điều chỉnh độ sáng (Dimmer), cảm biến nhiệt độ độ ẩm, công tắc thông minh nhiều kênh và ổ cắm thông minh.

\textbf{Ổn định hóa Occupancy Sensor:} Thay thế nút nhấn BTN0 bằng cảm biến chuyển động vật lý PIR hoặc cảm biến hiện diện radar mmWave thực tế, kết hợp lập trình các thuật toán chống nhiễu (debouncing/filtering) ở firmware để đưa ra trạng thái hiện diện chính xác nhất.

\textbf{Cấu hình nhóm (Zigbee Group/Scene):} Nghiên cứu triển khai tính năng nhóm thiết bị ở lớp mạng để thực hiện các lệnh điều khiển đồng thời (ví dụ: tắt toàn bộ đèn trong phòng) một cách đồng bộ và nhanh chóng.

\subsection{Nâng cấp tính năng cho Gateway nhúng}
\label{sec:ch6-gw-future}

\textbf{Chuyển đổi sang nền tảng máy chủ nhúng:} Chuyển tiến trình Z3Gateway từ máy tính xách tay Ubuntu sang chạy thực tế trên các máy tính nhúng như Raspberry Pi hoặc Orange Pi kết hợp với một module NCP EFR32 cắm qua cổng USB.

\textbf{Hoàn thiện Rules Engine cục bộ:} Phát triển bộ thông dịch rule động tại Gateway, cho phép người dùng cấu hình các kịch bản tự động hóa (ví dụ: Đèn tự bật khi Cảm biến báo có người vào ban đêm) ngay tại biên để hệ thống vẫn tự động hóa bình thường khi mất kết nối mạng.

\subsection{Nâng cấp nền tảng Cloud và Ứng dụng di động}
\label{sec:ch6-cloud-future}

\textbf{Truyền dữ liệu thời gian thực (Real-time update):} Bổ sung cầu nối WebSocket hoặc Server-Sent Events (SSE) giữa Cloud Backend và App di động để thay thế cho cơ chế kéo (poll) dữ liệu REST API hiện tại, giúp cập nhật trạng thái tức thời với băng thông tối thiểu.

\textbf{Triển khai tính năng OTA đã thiết kế:} Lập trình Gecko Bootloader trên thiết bị con và tích hợp logic nạp GBL file vô tuyến tại Gateway để hoàn thiện luồng nâng cấp phần mềm từ xa (OTA Campaign).

\textbf{Đẩy thông báo (Push Notification):} Tích hợp Firebase Cloud Messaging (FCM) để gửi thông báo khẩn cấp tới điện thoại người dùng khi phát hiện chuyển động bất thường hoặc khi có thiết bị ngoại tuyến.

\subsection{Tăng cường An ninh bảo mật}
\label{sec:ch6-security-future}

\textbf{Mã hóa TLS cho đường truyền MQTT:} Thiết lập chứng chỉ bảo mật SSL/TLS trên Mosquitto Broker và cấu hình Gateway/Cloud kết nối qua cổng 8883 để bảo mật toàn bộ dữ liệu trao đổi.

\textbf{Xác thực người dùng nâng cao:} Triển khai cơ chế xác thực OAuth2 kết hợp JWT (JSON Web Token) trên Cloud API để bảo vệ các endpoint điều khiển khỏi truy cập trái phép.

\newpage
